% Source: https://artofproblemsolving.com/community/c3h2023027_amcearly_aime_mock_problem
% Stable identity: aops-topic-2023027-q01
The condition says that the largest number minus the smallest is at most $2$. Count the ordered triples by this range.
\begin{itemize}
\item Range $0$: all three equal: $10$ triples.
\item Range $1$: the values are $k$ and $k+1$ for some $k=1,\dots,9$, and both values occur. Each person picks one of the two values, minus the $2$ constant choices: $2^3-2=6$ triples for each $k$, so $9\cdot6=54$.
\item Range $2$: the minimum is $k$ and the maximum is $k+2$ for some $k=1,\dots,8$. All three numbers lie in $\{k,k+1,k+2\}$ and both $k$ and $k+2$ occur; by inclusion-exclusion there are $3^3-2\cdot2^3+1=12$ such triples for each $k$, so $8\cdot12=96$.
\end{itemize}
The total is $10+54+96=\boxed{160}$. The answer $186$ suggested in the post is too large.
